\documentclass[12pt,a4paper]{article}

\usepackage[spanish,es-nodecimaldot]{babel}
\usepackage[utf8]{inputenc}
\usepackage[T1]{fontenc}
\usepackage{lmodern}
\usepackage{geometry}
\usepackage{microtype}
\usepackage{setspace}
\usepackage{parskip}
\usepackage{enumitem}
\usepackage{booktabs}
\usepackage{xcolor}
\usepackage{listings}
\usepackage{hyperref}
\usepackage{amsmath}
\usepackage{amssymb}

\geometry{margin=2.5cm}
\onehalfspacing
\setlist[itemize]{leftmargin=*,itemsep=0.2em}
\setlist[enumerate]{leftmargin=*,itemsep=0.2em}

\hypersetup{
  colorlinks=true,
  linkcolor=black,
  urlcolor=blue,
  pdftitle={Resumen teórico - Testing de Software (TP1)},
  pdfauthor={}
}

\lstdefinestyle{javastyle}{
  language=Java,
  basicstyle=\ttfamily\small,
  keywordstyle=\color{blue!60!black},
  commentstyle=\color{green!40!black},
  stringstyle=\color{orange!50!black},
  numbers=left,
  numberstyle=\tiny,
  stepnumber=1,
  numbersep=8pt,
  frame=single,
  breaklines=true,
  tabsize=2,
  showstringspaces=false
}

\title{\textbf{Resumen teórico}\\Testing de Software --- TP1}
\author{}
\date{}

\begin{document}

\maketitle
\tableofcontents
\newpage

\section*{Introducción}
\addcontentsline{toc}{section}{Introducción}
Este documento es un resumen teórico de las primeras unidades de \textit{Testing de Software} de la UNRC. Está pensado como material de estudio para acompañar el TP1 y se apoya en dos fuentes complementarias: los capítulos 1 y 2 del libro \textit{Introduction to Software Testing} (Ammann y Offutt, 2$^{a}$ edición) y las notas de clase de la cátedra (notas 00, 01 y 02 de la docente Valeria Bengolea).

La idea es presentar los conceptos en un hilo lógico que vaya desde el problema general --- por qué necesitamos testear software, qué relación tiene con la calidad y la confiabilidad --- hasta el aparato técnico que la materia va a usar a lo largo del cuatrimestre: la distinción entre \textit{fault}, \textit{error} y \textit{failure}, el modelo RIPR, la idea de criterios de cobertura y test requirements, los niveles y tipos de testing, y el proceso \textit{Model-Driven Test Design} (MDTD) que articula todo. Más que copiar definiciones, intentamos reorganizarlas, conectarlas entre sí y agregar las intuiciones que las hacen comprensibles.

\section{Calidad de software, confiabilidad y motivación del testing}
\subsection{Por qué testeamos software}
El testing existe porque construir software es una actividad de ingeniería sujeta a errores humanos, y el software que producimos termina sosteniendo aspectos críticos de la vida cotidiana: redes, energía, transporte, dispositivos médicos, sistemas embebidos, banca, comercio electrónico. Ammann y Offutt abren el capítulo 1 con esta observación: a diferencia de otras industrias, el software está en todas partes y, al mismo tiempo, casi siempre se asume que ``no va a fallar''. Esa asimetría entre la complejidad real y la expectativa optimista del usuario es lo que hace tan caro un defecto: cuando aparece, la consecuencia ya es visible para alguien.

Las notas refuerzan esta motivación con casos clásicos:
\begin{itemize}
  \item \textbf{Therac-25}: una condición de carrera en el control de un acelerador médico provocó dosis letales de radiación.
  \item \textbf{Ariane 5}: una conversión de punto flotante de 64 bits a entero de 16 bits, perfectamente correcta para Ariane 4, provocó la pérdida del cohete a los 37 segundos del lanzamiento. El defecto era ``de integración'': el software heredado nunca había sido reanalizado para la nueva trayectoria.
  \item \textbf{Mars Climate Orbiter}: un módulo trabajaba con unidades métricas y otro con unidades imperiales; la sonda se perdió al entrar a Marte.
  \item \textbf{Bug del Pentium (1994)}: una tabla mal inicializada en la circuitería de división de punto flotante. El análisis posterior mostró que casi cualquier criterio de cobertura unitario lo hubiera detectado.
  \item \textbf{Zune}: un bucle de cálculo del año actual con la regla de los años bisiestos se quedaba en \textit{loop} infinito durante el 31 de diciembre de un año bisiesto.
\end{itemize}

El reporte del NIST de 2002 citado tanto en el libro como en las notas pone una cifra concreta: un testing inadecuado cuesta a la economía estadounidense entre 22.000 y 59.000 millones de dólares anuales, y mejoras razonables en QA podrían reducir esa cifra a la mitad. Las fallas en aplicaciones financieras se miden directamente en millones por hora. La conclusión es directa: el testing no es un lujo de fin de proyecto, es una inversión.

\subsection{Calidad, confiabilidad y la naturaleza particular de las fallas de software}
La calidad del software se define a través de varios atributos: confiabilidad, mantenibilidad, interoperabilidad, seguridad, eficiencia, usabilidad, entre otros. De todos ellos la \textbf{confiabilidad} es uno de los más importantes y el más vinculado al testing. Una medida tradicional, aunque limitada, es la cantidad de defectos por línea de código: a mayor densidad de defectos, menos confiable es el producto.

Hay un punto sutil pero fundamental: las fallas de software \textbf{no son como las fallas mecánicas o eléctricas}. Una pieza mecánica falla por desgaste; un programa, no. El defecto que produce una falla de software está presente desde el momento mismo en que se escribió el código (o se diseñó el sistema, o se redactó el requisito) y solo se manifiesta cuando un escenario particular lo activa. Esto cambia la forma de pensar: no buscamos prevenir un deterioro, buscamos descubrir un error de diseño que ya está ahí.

\subsection{Verificación y Validación (V\&V)}
La actividad general que abarca al testing se denomina \textbf{V\&V}, y conviene distinguir sus dos componentes:

\begin{itemize}
  \item \textbf{Verificación}: ``¿estamos construyendo el producto correctamente?''. Es decir, ¿el software cumple con la especificación que tenemos escrita?
  \item \textbf{Validación}: ``¿estamos construyendo el producto correcto?''. Es decir, ¿la especificación misma refleja lo que el usuario realmente necesita?
\end{itemize}

Las definiciones formales del libro lo dicen así (Definiciones 1.4 y 1.5): verificación es el proceso de determinar si los productos de una fase del desarrollo satisfacen los requisitos establecidos en la fase previa; validación es el proceso de evaluar el software al final del desarrollo para asegurar que cumple con el uso pretendido. La verificación es más técnica y se apoya en especificaciones y requisitos formales; la validación depende del conocimiento del dominio.

Una intuición útil que aparece en las notas es ver el flujo
\[
\textit{requisitos reales} \;\longrightarrow\; \textit{especificación} \;\longrightarrow\; \textit{implementación}
\]
como dos saltos distintos: la validación atraviesa todo el flujo y pregunta si el resultado refleja los requisitos reales; la verificación se concentra en el último salto y pregunta si la implementación cumple con la especificación.

El problema general de la verificación se suele escribir de manera muy compacta como
\[
\textit{Sys} \models \textit{Spec}
\]
es decir: el sistema se comporta de acuerdo a su especificación. Operacionalmente, para programas imperativos, una manera muy útil de expresar esa especificación es mediante aserciones de corrección $\{P\}\, a\, \{Q\}$: siempre que el programa $a$ comience en un estado que satisface la precondición $P$, debe terminar en un estado que satisface la postcondición $Q$.

\subsection{Especificaciones, contratos e invariantes}
Las especificaciones son fundamentales para razonar sobre la corrección y, por lo tanto, para diseñar tests con criterio. En las clases se trabajan tres niveles de formalidad:
\begin{itemize}
  \item \textbf{Informales}: comentarios en lenguaje natural.
  \item \textbf{Semi-informales}: tags estilo Javadoc con \texttt{@precondition} y \texttt{@postcondition} escritas en prosa controlada.
  \item \textbf{Formales}: lenguajes como JML, que producen especificaciones \textit{chequeables} mecánicamente.
\end{itemize}

Las precondiciones y postcondiciones establecen \textbf{contratos} entre cliente y proveedor con derechos y obligaciones simétricas:
\begin{itemize}
  \item El \textbf{cliente} \textit{debe} garantizar la precondición antes de llamar a la rutina y \textit{puede} suponer la postcondición al regreso.
  \item El \textbf{proveedor} \textit{puede} suponer la precondición al recibir el control y \textit{debe} garantizar la postcondición al terminar.
\end{itemize}

Cuando una aserción se viola en tiempo de ejecución, el responsable depende de cuál se violó: si fue la precondición, el error es del cliente; si fue la postcondición, del proveedor. Esto convierte a las aserciones en un mecanismo de \textit{trazabilidad} del defecto.

A nivel de clases, además de pre/postcondiciones de cada método, se utilizan \textbf{invariantes de representación} (también llamados invariantes de clase): propiedades globales que toda instancia debe satisfacer. Las reglas son sencillas:
\begin{itemize}
  \item Todos los constructores deben \textbf{establecer} el invariante al finalizar.
  \item Todos los métodos deben \textbf{preservarlo}: si el invariante valía antes de la llamada, debe seguir valiendo después.
\end{itemize}

Formalmente, esto se expresa pidiendo que cada método $m$ satisfaga la tripla $\{\text{pre-}m \wedge \text{Inv}\}\; m\; \{\text{post-}m \wedge \text{Inv}\}$. En código, el invariante suele codificarse como un método \texttt{repOk()} que devuelve \texttt{true} si y solo si el objeto está bien formado; es una herramienta muy útil para testing y debugging.

\section{Análisis estático y análisis dinámico}
El proceso de V\&V tiene dos grandes objetivos: \textit{descubrir defectos} y \textit{medir} si el sistema es usable en situaciones reales. Para lograrlo se combinan técnicas de \textbf{análisis} de los artefactos del software (programas, modelos, especificaciones, documentos). Para el código en particular, esas técnicas se agrupan en dos familias.

\subsection{Análisis dinámico}
Es el análisis que infiere propiedades del programa \textbf{ejecutándolo}. Tiene dos cualidades complementarias:
\begin{itemize}
  \item Es \textbf{preciso}: lo que se observa efectivamente ocurrió en la corrida.
  \item Es \textbf{incompleto}: no cubre todos los comportamientos posibles, solo los inducidos por las entradas usadas.
\end{itemize}

Bajo este paraguas entran el \textit{testing}, el \textit{profiling} (consumo de recursos como memoria) y el chequeo en tiempo de ejecución de propiedades del programa.

\subsection{Análisis estático}
Infiere propiedades a partir del \textbf{texto} del programa, sin ejecutarlo. Es complementario al dinámico: en general puede cubrir todos los caminos posibles, pero a costa de ser \textbf{impreciso} (suele dar falsos positivos). Incluye inspección manual, chequeo de tipos, verificación estática de programas y herramientas de análisis. Una inspección manual cuidadosa puede ser sorprendentemente efectiva para encontrar defectos, aunque rara vez es exhaustiva.

\subsection{Una observación clásica de Dijkstra}
La frase más citada de la materia pertenece a Dijkstra (1970):
\begin{quote}
``Program testing can be used to show the presence of bugs, but never to show their absence.''
\end{quote}
Esta cita captura el límite teórico del testing: aunque ejecutemos millones de tests, nunca probamos formalmente que no haya defectos. Lo que sí podemos hacer es \textit{reducir el riesgo} y \textit{aumentar la confianza}, y para eso necesitamos un proceso planificado, criterios objetivos y herramientas que cubran sistemáticamente el espacio de comportamientos del programa.

\section{Terminología fundamental: fault, error, failure, bug}
Una de las herramientas más útiles del curso es el vocabulario preciso para hablar de defectos. Conviene fijarlo desde el principio porque, en el habla cotidiana, ``bug'' se usa de manera ambigua para tres cosas distintas.

\subsection{Definiciones (capítulo 1, Ammann \& Offutt)}
\begin{itemize}
  \item \textbf{Fault (defecto)}: ``A static defect in the software''. Es el problema estático que vive en el código, el diseño o la especificación. Existe se ejecute o no el programa.
  \item \textbf{Error}: ``An incorrect internal state that is the manifestation of some fault''. Es un estado interno incorrecto que aparece cuando el defecto se activa durante la ejecución.
  \item \textbf{Failure (falla)}: ``External, incorrect behavior with respect to the requirements or another description of the expected behavior''. Es la manifestación externa observable del problema.
\end{itemize}

La cadena causal habitual es
\[
\textit{fault} \;\longrightarrow\; \textit{error} \;\longrightarrow\; \textit{failure}.
\]

Un \textbf{bug} es un término informal que en el lenguaje cotidiano agrupa a los tres. En el contexto de la materia preferimos los términos específicos.

\subsection{Analogía médica}
Para fijar las ideas, el libro y las notas usan una analogía con la medicina: un paciente llega al médico con una lista de \textit{síntomas} (esos son las \textit{failures}); el médico intenta diagnosticar la causa (el \textit{fault}); para hacerlo puede medir condiciones anómalas internas como presión, arritmias o niveles de glucosa (esos son los \textit{errors}). La analogía aclara la jerarquía, pero hay una diferencia crucial: los fallos de software no provienen de un ``desgaste natural'' como muchas enfermedades, sino de \textbf{errores de diseño} humanos que estuvieron presentes desde el comienzo.

\subsection{Ejemplo concreto: \texttt{numZero}}
El libro presenta un método Java muy simple para ilustrar la diferencia entre los tres conceptos:
\begin{lstlisting}[style=javastyle]
public static int numZero(int[] x) {
    int count = 0;
    for (int i = 1; i < x.length; i++) {  // i = 1 está mal
        if (x[i] == 0) count++;
    }
    return count;
}
\end{lstlisting}
El defecto (\textit{fault}) es comenzar el recorrido en \texttt{i = 1} en lugar de \texttt{i = 0}.

\begin{itemize}
  \item Con \texttt{x = [2, 7, 0]} el método devuelve \texttt{1}, que es el resultado correcto. El defecto se ejecuta y produce un \textbf{error} (el valor de \texttt{i} es incorrecto en la primera iteración), pero el error \textit{no propaga} a la salida y \textbf{no hay falla observable}.
  \item Con \texttt{x = [0, 7, 2]} el método devuelve \texttt{0} en lugar de \texttt{1}. Acá el error sí propaga a la variable \texttt{count}, y se traduce en una \textbf{falla} visible.
\end{itemize}

Este ejemplo muestra algo crucial: \textbf{ejecutar el defecto no garantiza ver la falla}. Para que la falla se manifieste se tienen que dar varias condiciones, que el modelo RIPR formaliza.

\section{El modelo RIPR}
RIPR --- por \textbf{Reachability}, \textbf{Infection}, \textbf{Propagation}, \textbf{Revealability} --- es el modelo central que la materia usa para razonar sobre cuándo y por qué un test revela un defecto. Aparece en el capítulo 2 del libro y se discute en la nota 02 con un diagrama que ilustra el flujo.

\subsection{Las cuatro condiciones}
Para que un test logre revelar un defecto deben darse, en orden, cuatro condiciones:

\begin{enumerate}
  \item \textbf{Reachability (alcanzar)}: la entrada del test tiene que llevar la ejecución hasta el lugar (o lugares) del programa donde está el defecto.
  \item \textbf{Infection (infectar)}: una vez ejecutado el código defectuoso, el estado del programa tiene que volverse incorrecto.
  \item \textbf{Propagation (propagar)}: el estado infectado tiene que seguir vivo y propagarse a través del resto de la ejecución hasta provocar que alguna salida o estado final sea incorrecto.
  \item \textbf{Revealability (revelar)}: el tester tiene que efectivamente \textit{observar} la porción incorrecta del estado final, comparándola contra un oráculo. Si solo observa partes correctas, la falla queda oculta.
\end{enumerate}

\subsection{Por qué importa cada eslabón}
RIPR explica por qué tantos defectos sobreviven al testing aunque parezca exhaustivo:
\begin{itemize}
  \item Si no hay \textit{reachability}, el código defectuoso nunca se ejecuta (por ejemplo, una condición que solo se da en producción).
  \item Si no hay \textit{infection}, el código se ejecuta pero produce el mismo valor que la versión correcta para esa entrada (caso \texttt{numZero([2,7,0])}).
  \item Si no hay \textit{propagation}, el estado interno está mal pero se sobrescribe antes de salir.
  \item Si no hay \textit{revealability}, la falla está en la salida pero el oráculo no la mira.
\end{itemize}

Este modelo aplica incluso cuando el defecto es por \textit{omisión}: cuando falta código que debería estar, el contador de programa al pasar por esa zona ``ausente'' ya constituye un estado erróneo. Diseñar tests con RIPR en mente es lo que diferencia un test que ``llama al método'' de uno que efectivamente busca exponer un defecto.

\section{Testing y debugging: dos procesos distintos}
Es habitual confundir \textit{testing} con \textit{debugging}; las notas y el libro los separan con cuidado.

\begin{itemize}
  \item \textbf{Testing}: ``Evaluating software by observing its execution''. Es el proceso de \textit{evaluar} el comportamiento del software ejecutándolo. El resultado de un test es ``pasa / falla''.
  \item \textbf{Test Failure}: ejecución de un test que da lugar a una falla en el software.
  \item \textbf{Debugging}: ``The process of finding a fault given a failure''. Es lo que viene \textit{después} de un test que falla: a partir de la falla observada, buscamos el defecto en el código que la causó.
\end{itemize}

Un test puede ser \textit{exitoso} (no revela falla) o \textit{fallido} (descubre una falla). En testing, paradójicamente, un test ``fallido'' es \textit{información valiosa}: nos dice dónde mirar.

\subsection{Definición de testing por niveles de pensamiento}
Beizer propuso una caracterización del testing en cinco \textit{niveles de madurez} según cómo lo entiende cada profesional. El libro y las notas la adoptan:

\begin{itemize}
  \item \textbf{Nivel 0}: no hay diferencia entre testing y debugging. Es la postura ingenua del estudiante de programación: ``hago compilar, pruebo con un par de entradas, corrijo lo que sale mal''.
  \item \textbf{Nivel 1}: el propósito del testing es mostrar \textbf{corrección}. Es un paso adelante, pero teóricamente imposible: la corrección no se puede establecer ejecutando un número finito de tests.
  \item \textbf{Nivel 2}: el propósito del testing es mostrar que \textbf{el software no funciona}. Es un objetivo válido (encontrar fallas), pero pone a testers y desarrolladores en relación adversaria.
  \item \textbf{Nivel 3}: el propósito no es probar nada específico, sino \textbf{reducir el riesgo} asociado al uso del software. Acepta que toda decisión de usar software implica algún riesgo y orienta el trabajo a controlarlo.
  \item \textbf{Nivel 4}: testing es una \textbf{disciplina mental} que ayuda a los profesionales a desarrollar software de mejor calidad. La responsabilidad principal no es solo medir, sino \textit{mejorar} la calidad.
\end{itemize}

Esta escala es útil porque obliga a preguntarse en qué nivel está la organización (o uno mismo) y qué cambios harían falta para pasar al siguiente.

\section{Casos de prueba y cómo definirlos}
\subsection{Componentes de un test}
Un test ---para ser útil--- no es solo una entrada. Las notas lo descomponen en tres partes mínimas:
\begin{enumerate}
  \item \textbf{Entradas}: los valores que se pasan al código bajo prueba.
  \item \textbf{Código a ejecutar (SUT, \textit{Software Under Test})}: el método o componente que se testea.
  \item \textbf{Descripción de los resultados esperados (condición de aceptación)}: el \textit{oráculo}, contra el cual se contrasta lo que realmente devolvió el programa.
\end{enumerate}

En la práctica de JUnit esto se materializa así:
\begin{lstlisting}[style=javastyle]
@Test
void oddOrPosTest() {
    int[] x = {-3, -2, 0, 1, 4};
    int result = SimpleRoutines.oddOrPos(x);
    assertEquals(3, result, "Array contains positive number");
}
\end{lstlisting}
La elección de la entrada \texttt{\{-3,-2,0,1,4\}} no es casual ---tiene un impar negativo, un par negativo, un cero, un impar positivo y un par positivo--- y la aserción \texttt{assertEquals(3, result)} actúa como oráculo.

\subsection{Buenas prácticas al definir tests}
Las notas (siguiendo a Myers y otros) listan varios principios:
\begin{itemize}
  \item Un programador debería evitar testear sus propios programas: psicológicamente no estamos preparados para buscar errores en nuestro propio trabajo, y un error frecuente es que el programador no entendió bien el enunciado o la especificación, en cuyo caso testeará el mismo malentendido.
  \item Esto está cambiando con metodologías ágiles y \textit{Test-Driven Development}, donde escribir tests es parte natural del trabajo del desarrollador.
  \item Probar que el software \textit{no haga lo que debe hacer} es solo la mitad del trabajo: la otra mitad es probar que \textit{no haga lo que no debe hacer}.
  \item Evitar los \textbf{tests desechables}: si tuviéramos que reinventar los tests cada vez que volvemos al programa, perdemos todo el trabajo. Guardar tests y volver a ejecutarlos tras cambios en otras componentes es lo que se llama \textbf{regression testing} (testing de regresión).
\end{itemize}

\subsection{Las dos preguntas centrales al diseñar tests}
Las notas formulan dos preguntas que vuelven a aparecer una y otra vez en la materia:
\begin{itemize}
  \item ¿Cómo \textbf{elegimos} las entradas?
  \item ¿Contra qué \textbf{contrastamos} el resultado obtenido (oráculo)?
\end{itemize}
El resto de la materia se dedica a dar respuestas progresivamente más sofisticadas a estas dos preguntas.

\section{Procesos de desarrollo y rol del testing}
\subsection{Procesos clásicos}
Los modelos clásicos de proceso de desarrollo (estilo cascada) involucran fases relativamente bien separadas: requisitos, especificación, planeamiento, diseño, implementación, integración y testing, mantenimiento. En estos modelos, el testing aparece principalmente como una fase \textit{posterior a la implementación}.

\subsection{Procesos ágiles}
En los modelos ágiles el testing está presente en varias fases del proceso, principalmente como parte de la codificación misma. Cada nueva funcionalidad se acompaña de tests; el ciclo típico es escribir código $\to$ mejorar tests $\to$ ejecutar tests $\to$ reparar si hay errores $\to$ repetir hasta cubrir la especificación.

\subsection{Test-Driven Development (TDD)}
En TDD se da una vuelta de tuerca adicional: los \textbf{tests se escriben antes de la codificación} y funcionan como \textit{especificación de la funcionalidad} y \textit{criterio de aceptación} del código. El ciclo es: escribir tests $\to$ escribir código para que pasen $\to$ ejecutar tests $\to$ reparar / mejorar $\to$ repetir.

\subsection{Costo de detectar defectos tarde}
Una idea muy importante del capítulo 1 (con la famosa Figura 1.1) es que el \textbf{costo} de encontrar y reparar un defecto crece dramáticamente cuanto más tarde se lo detecta:
\begin{itemize}
  \item Detectarlo durante requisitos, diseño o testing unitario: costo unitario.
  \item Durante integración: $\sim 5\times$.
  \item Durante system testing: $\sim 10\times$.
  \item Después del despliegue: $\sim 50\times$.
\end{itemize}
Por eso la consigna estratégica de Beizer (nivel 4) se condensa en: ``eliminar defectos tan temprano como sea posible''. Cada falla atrapada en testing unitario ahorra dinero.

\section{Niveles de testing}
Los niveles de testing son una categorización clásica, frecuentemente representada con el \textbf{modelo en V}, que asocia a cada fase del desarrollo un nivel de testing.

\subsection{Niveles tradicionales (modelo en V)}
\begin{itemize}
  \item \textbf{Acceptance testing}: evalúa el software respecto a los requisitos y necesidades del usuario. ¿Hace el software lo que el cliente quería?
  \item \textbf{System testing}: evalúa el sistema completo contra su diseño arquitectónico y su comportamiento global. Suele encontrar fallas costosas porque ya hay integración entre módulos.
  \item \textbf{Integration testing}: evalúa cómo se comunican entre sí los módulos. Asume que cada módulo funciona bien individualmente y se concentra en las \textit{interfaces}.
  \item \textbf{Module testing}: evalúa cada módulo (típicamente una clase o un paquete) en aislamiento, incluyendo cómo interactúan internamente sus unidades.
  \item \textbf{Unit testing}: evalúa la unidad más chica (un método o procedimiento) de manera aislada. Es el nivel ``más bajo''.
\end{itemize}

\subsection{Testing de aceptación: usuario vs. negocio}
El testing de aceptación se subdivide a su vez en dos:
\begin{itemize}
  \item \textbf{Aceptación de usuario}: lo realiza el usuario para determinar si el sistema cubre sus necesidades. Es esencialmente \textit{validación}.
  \item \textbf{Aceptación de negocio}: lo realiza el equipo de desarrollo/testing para determinar si el sistema eventualmente pasará las pruebas del usuario. Es esencialmente \textit{verificación}.
\end{itemize}

\subsection{Niveles específicos de programación orientada a objetos}
El testing en orientación a objetos rompe un poco la dicotomía unit/module, porque las clases agrupan muchos métodos relacionados. Por eso se introducen niveles más finos:
\begin{itemize}
  \item \textbf{Intra-method testing}: testear cada método individualmente.
  \item \textbf{Inter-method testing}: testear pares de métodos de la misma clase (cómo interactúan vía el estado del objeto).
  \item \textbf{Intra-class testing}: testear una clase entera mediante secuencias de llamadas a sus métodos.
  \item \textbf{Inter-class testing}: testear varias clases juntas; es una forma de testing de integración.
\end{itemize}

\subsection{Tipos básicos de testing por su objetivo}
Independientemente del nivel, las notas mencionan algunos tipos de testing que se distinguen por su \textit{objetivo} o \textit{condición de aceptación}:
\begin{itemize}
  \item \textbf{Regression}: la condición de aceptación es que se preserve el comportamiento de versiones anteriores del software. Se guardan los tests y se vuelven a ejecutar tras cada cambio.
  \item \textbf{Diferencial}: se contrasta el comportamiento contra otro software con el mismo propósito.
  \item \textbf{Aceptación}: se evalúa contra los requisitos o necesidades del usuario.
\end{itemize}

\section{Criterios de cobertura y test requirements}
\subsection{El problema del espacio de entradas}
Aunque parezca que un método como
\begin{lstlisting}[style=javastyle]
private static double computeAverage(int A, int B, int C)
\end{lstlisting}
es trivial, el espacio de entradas posibles es gigantesco: tres enteros de 32 bits dan $\sim 80 \cdot 10^{27}$ combinaciones distintas. A efectos prácticos, el espacio de inputs es \textbf{infinito}. Ejecutar ``todos los tests'' es inviable.

El tester se enfrenta entonces a dos preguntas estructurales:
\begin{enumerate}
  \item ¿Cómo \textbf{buscamos} en ese espacio?
  \item ¿Cuándo \textbf{paramos} de buscar?
\end{enumerate}

\subsection{Qué es un criterio de cobertura}
Los \textbf{criterios de cobertura} son la herramienta principal para responder esas dos preguntas. Un criterio de cobertura es una regla (o un conjunto de reglas) que define un conjunto de \textit{test requirements}: elementos específicos de un artefacto que cada test debe satisfacer o cubrir.

\begin{itemize}
  \item \textbf{Criterio de test}: colección de reglas y un proceso que define los test requirements. Ejemplos: ``cubrir cada línea de código'', ``cubrir cada requisito funcional''.
  \item \textbf{Test requirement}: elemento específico de un artefacto de software que un test debe satisfacer. Cada línea de código (o cada rama, o cada par de aristas\ldots) puede dar lugar a un test requirement.
\end{itemize}

Satisfacer un criterio de cobertura le da al tester dos garantías importantes:
\begin{itemize}
  \item La búsqueda en el espacio de entradas se hizo de manera \textbf{minuciosa}: no quedan rincones obvios sin explorar.
  \item Los tests no tienen \textbf{demasiado solapamiento}: no estamos redundando trabajo.
\end{itemize}

Además, los criterios proveen una \textbf{regla de parada} (\textit{stopping rule}): se sabe \textit{a priori} cuántos tests son necesarios para satisfacer el criterio y se sabe cuándo se terminó.

\subsection{Cuatro estructuras matemáticas, muchos criterios}
Una observación profunda del libro de Ammann \& Offutt, que organiza toda la segunda parte de la materia, es que los cientos de criterios de cobertura definidos en la literatura se pueden \textbf{reducir} a unos pocos sobre cuatro tipos de estructuras matemáticas:
\begin{enumerate}
  \item \textbf{Dominio de entradas} (\textit{input domain}): particiones del espacio de entradas.
  \item \textbf{Grafos}: grafos de control de flujo, grafos de llamadas, máquinas de estado.
  \item \textbf{Expresiones lógicas}: condiciones, decisiones, MC/DC.
  \item \textbf{Descripciones sintácticas}: gramáticas, mutaciones del programa.
\end{enumerate}

El criterio ``cobertura de pares de aristas'' en un grafo de control de flujo, por ejemplo, generará un test requirement por cada par de aristas consecutivas; el tester deberá encontrar caminos (y luego inputs concretos) que recorran todos esos pares.

\section{Black-box, white-box y la perspectiva de MDTD}
Una clasificación tradicional del testing es:

\begin{itemize}
  \item \textbf{Testing de caja negra (black-box)}: los tests se generan a partir de descripciones \textit{externas} del software --- especificaciones, requisitos, diseño. No se mira el código.
  \item \textbf{Testing de caja blanca (white-box)}: los tests se generan a partir del \textit{código} mismo, específicamente sus ramas, condiciones y líneas. Algunos autores agregan ``gray-box'' para tests que parten del diseño interno.
\end{itemize}

Hay una observación importante del libro: aunque históricamente se asoció black-box con testing de sistema y white-box con unit testing, esa correspondencia es \textit{artificial}. Cualquiera de los dos enfoques se puede aplicar a cualquier nivel.

Más aún, desde la perspectiva de \textbf{MDTD} (Model-Driven Test Design) la distinción pierde importancia: lo que importa es \textbf{de qué nivel de abstracción se deriva el modelo} con el que se diseña el test. Un grafo de control de flujo extraído del código y un grafo extraído de un diagrama UML se manejan con los mismos criterios; lo que cambia es la fuente.

\section{Actividades del testing y Model-Driven Test Design (MDTD)}
\subsection{Las cuatro actividades centrales}
El capítulo 2 separa el trabajo de testing en cuatro actividades discretas, cada una con habilidades, conocimientos y formación distintos:

\begin{enumerate}
  \item \textbf{Diseño de tests (test design)}: diseñar valores de entrada que efectivamente ejerciten el software. Se subdivide en dos sub-enfoques:
    \begin{itemize}
      \item \textbf{Basado en criterios}: diseñar tests para satisfacer un criterio de cobertura u otro objetivo formal. Es la parte más matemática del testing.
      \item \textbf{Basado en factor humano} (\textit{human-based} o \textit{stress tests}): diseñar tests aprovechando el conocimiento del dominio y del comportamiento típico/extremo de los usuarios. Busca activamente entradas raras, valores límite, situaciones no previstas.
    \end{itemize}
  \item \textbf{Automatización de tests (test automation)}: codificar los tests en \textit{scripts ejecutables} (por ejemplo, casos de JUnit). Requiere habilidad de programación.
  \item \textbf{Ejecución de tests (test execution)}: correr los tests sobre el software y registrar los resultados. Si están bien automatizados es trivial.
  \item \textbf{Evaluación de tests (test evaluation)}: evaluar los resultados, interpretarlos y reportarlos a desarrolladores o gerentes. Es mucho más difícil de lo que parece porque requiere conocimiento del dominio, de testing y de comunicación.
\end{enumerate}

Además de estas cuatro, hay actividades de soporte: \textbf{gestión de tests} (organización del equipo, elección de criterios, grado de automatización), \textbf{mantenimiento de tests} (guardar y actualizar tests a medida que el software evoluciona) y \textbf{documentación de tests} (criterios y requisitos cubiertos, justificación de tests diseñados a mano).

Los enfoques basados en criterios y basados en humano \textbf{no compiten}, se complementan: lo cuantitativo (criterios formales) y lo cualitativo (intuición humana del dominio) cubren tipos de defectos distintos.

\subsection{El proceso MDTD}
MDTD propone un proceso explícito de diseño de tests que ``eleva el nivel de abstracción''. La idea es exactamente la misma que en ingeniería tradicional: en lugar de razonar directamente sobre la implementación, se construye un \textbf{modelo matemático} (un grafo, una partición, una expresión lógica, una gramática), se trabaja a ese nivel y luego se baja a valores concretos.

El flujo, ilustrado en la figura 2.4 del libro y en las transparencias de la nota 02, se puede leer como una cadena de transformaciones:

\begin{enumerate}
  \item \textbf{Artefacto software} (código, UML, requisitos, manual de usuario).
  \item De ese artefacto se extrae un \textbf{Modelo / Estructura} (por ejemplo, un grafo de control de flujo).
  \item Aplicando un criterio sobre ese modelo se generan \textbf{Test requirements}.
  \item Los test requirements se refinan a \textbf{especificaciones de tests} (caminos concretos a recorrer, valores a probar).
  \item Se eligen \textbf{Valores de inputs} que satisfacen las especificaciones.
  \item Esos valores se empaquetan como \textbf{tests}.
  \item Los tests se traducen a \textbf{scripts ejecutables}.
  \item Se obtienen \textbf{resultados} y, finalmente,
  \item se evalúa cada resultado como \textbf{pasa / falla}.
\end{enumerate}

Lo de arriba se divide en dos planos: el plano de \textit{diseño} (artefacto $\to$ modelo $\to$ requisitos $\to$ especificaciones) y el plano de \textit{implementación} (valores $\to$ tests $\to$ scripts $\to$ resultados $\to$ evaluación). MDTD separa también las cuatro actividades del punto anterior: el diseño, la automatización, la ejecución y la evaluación quedan claramente delimitadas y pueden ser realizadas por personas distintas.

\subsection{Ejemplo mínimo de MDTD}
El propio capítulo 2 ilustra el proceso con un método Java muy simple:
\begin{lstlisting}[style=javastyle]
/** Return index of node n at the first position it appears,
 *  -1 if it is not present */
public int indexOf(Node n) {
    for (int i = 0; i < path.size(); i++)
        if (path.get(i).equals(n))
            return i;
    return -1;
}
\end{lstlisting}
A partir del código se extrae un grafo de control de flujo con cinco nodos (inicio, guarda del \texttt{for}, evaluación del \texttt{if}, \texttt{return i}, \texttt{return -1}) y aristas que reflejan el flujo. Si se elige el criterio de \textbf{cobertura de pares de aristas}, se obtienen seis test requirements (subcaminos como \texttt{[1,2,3]}, \texttt{[1,2,5]}, \texttt{[2,3,4]}, etc.), y se buscan tres caminos completos que entre todos cubran los seis pares. Recién al final se eligen valores concretos de \texttt{path} y \texttt{n} que recorran esos caminos.

\section{Por qué testeamos: síntesis de objetivos}
A partir de todo lo anterior, las notas cierran la unidad de introducción con una lista compacta de \textbf{objetivos} del testing:

\begin{itemize}
  \item Revelar fallas tan temprano como sea posible durante el desarrollo.
  \item \textbf{Mejorar la calidad} del software: si encontramos fallas, las corregimos antes del despliegue.
  \item \textbf{Reducir el costo} de mantenimiento y operación: cada falla atrapada en testing unitario evita costos $50\times$ mayores tras despliegue.
  \item Mantener a los \textbf{usuarios satisfechos} y proteger la confianza en el producto.
  \item Más estratégicamente (nivel 3), \textbf{reducir el riesgo} asociado al uso del software.
  \item En el nivel más maduro (nivel 4), \textbf{aumentar la calidad} de manera sostenida formando a los desarrolladores: medir, reportar, formar, prevenir.
\end{itemize}

\section{Síntesis final: ideas clave para fijar}
Como cierre, las ideas más importantes que deberían quedar fijas después de leer este resumen son las siguientes:

\begin{enumerate}
  \item \textbf{Testing no es debugging}. Testing evalúa el software para revelar fallas; debugging busca el defecto causal a partir de una falla observada. Son procesos distintos con habilidades distintas.
  \item \textbf{Fault, error, failure} forman una cadena causal: el defecto vive en el código, se manifiesta como un estado interno incorrecto cuando se ejecuta, y produce una falla cuando ese estado infectado se propaga a una salida observable.
  \item \textbf{El modelo RIPR (Reachability, Infection, Propagation, Revealability)} es el marco para diseñar tests que realmente expongan defectos: alcanzar la línea, infectar el estado, propagar el error a la salida y observarlo con un oráculo.
  \item \textbf{El testing no prueba ausencia de defectos}, solo presencia (Dijkstra). Por eso se trabaja con criterios objetivos que sistematicen la exploración y provean reglas de parada.
  \item \textbf{Verificación y validación} son distintas: la primera mira si construimos el producto correctamente respecto a la especificación; la segunda si construimos el producto correcto respecto a las necesidades reales.
  \item Las \textbf{especificaciones} ---informales, semi-informales o formales (JML)--- son la base para razonar sobre corrección. Los \textbf{contratos} de pre/postcondiciones y los \textbf{invariantes de clase} estructuran ese razonamiento y permiten asignar responsabilidad ante cada violación.
  \item Un \textbf{test} se compone de entradas, código a ejecutar y oráculo. Sin oráculo el test no sirve: nunca sabremos si lo que pasó está bien o mal.
  \item Los \textbf{criterios de cobertura} son la respuesta estructurada al problema del espacio infinito de entradas: definen un conjunto finito de \textit{test requirements} que organiza la búsqueda y dicta cuándo parar.
  \item Todo criterio razonable se puede definir sobre \textbf{cuatro estructuras matemáticas}: dominio de entradas, grafos, expresiones lógicas y descripciones sintácticas. Esta es la columna vertebral de la materia.
  \item Los \textbf{niveles de testing} (unit, module, integration, system, acceptance) y sus variantes para POO (intra-method, inter-method, intra-class, inter-class) describen \textit{a qué nivel del software apunta el test}, no \textit{cómo se diseña}.
  \item La \textbf{automatización} (JUnit, scripts) es lo que vuelve viable el regression testing y, en última instancia, la capacidad de mantener un software vivo sin romperlo.
  \item \textbf{MDTD (Model-Driven Test Design)} es el proceso global: del artefacto al modelo, del modelo a los requisitos, de los requisitos a valores concretos, de los valores a scripts y al veredicto pasa/falla. Separa cuatro actividades (diseño, automatización, ejecución, evaluación) que requieren perfiles distintos.
  \item El testing es, en última instancia, una \textbf{disciplina de calidad}: detectar defectos temprano abarata, reduce el riesgo del uso del software y, en organizaciones maduras, mejora a quienes lo construyen.
\end{enumerate}

Con este marco como base, las próximas unidades (cobertura de grafos, cobertura lógica, particionado del espacio de entradas, testing basado en sintaxis, property-based testing y generación automática) van a verse como aplicaciones concretas de las mismas ideas: un modelo, un criterio que produce test requirements, valores que los satisfacen y un oráculo que decide pasa/falla.

\end{document}
